\section{Virtual Lab 实战：COVID 变体 nanobody 设计案例}

课程中最具说服力的部分是具体案例：让 Virtual Lab 针对 SARS-CoV-2 新变体设计 binder。这个任务有真实约束：时间紧、数据分布变动快、可用公开数据不完备，且最终需要实验验证而非停留在文本推理。

\subsection{从``做抗体''到``做 nanobody''的路径转移}

在早期讨论中，immunology 角色提出了较反直觉建议：优先设计 nanobody 而不是常规 antibody。Zou 强调这是一个``如果问多数人类研究者，未必先走这条路''的决策点。其后 machine learning 角色给出可计算性理由：nanobody 更小，结构建模与打分稳定性更高，更匹配现有计算工具链。

这段讨论体现了 agentic team 的价值：建议不是凭``灵感''，而是由跨角色互证形成。immunology 给生物学可行性，ML 给计算约束与可执行性，再由 PI 进行任务收敛。

\begin{importantbox}{案例中的关键决策逻辑}
\begin{itemize}
    \item 目标不是追求理论上最优 binder，而是追求\textbf{在当前工具条件下更可验证}的方案。
    \item 路线选择同时受生物机理与计算可操作性约束。
    \item 系统先选择``可持续迭代''路径，而不是一次性豪赌。
\end{itemize}
\end{importantbox}

\subsection{Critic Agent 的作用：防止高置信错误累积}

在该案例中，Critic 明确指出：nanobody 公共数据规模较小，若直接按常规深度学习流程推进，容易出现过拟合并导致虚假高分。这个提醒并没有否定主路线，而是迫使团队增加校验与保守性假设。

图\ref{fig:vl-discussion} 对应课程中的讨论总结页。\footnote{来源：Berkeley RDI 课程视频，关键帧时间戳 00:22:10。}

\begin{figure}[H]
\centering
\includegraphics[width=0.9\textwidth]{frame-02-agent-meeting.jpg}
\caption{Virtual Lab 讨论页：多 Agent 协作优于单 Agent，并强调 memory/sandbox}
\label{fig:vl-discussion}
\end{figure}

\begin{knowledgebox}{读图：Discussion 页是系统需求而非宣传结论}
四个要点分别对应任务覆盖、专家互补、执行隔离和多体行为。memory 保存跨轮证据，sandbox 限制代码与工具副作用，social dynamics 则提醒协作可能产生从众、议题漂移和地位偏置。因此“多 Agent 优于单 Agent”只在协议、异质性和审查机制同时成立时才有意义。
\end{knowledgebox}

\begin{warningbox}{没有 Critic 的多 Agent 系统会怎样}
\begin{itemize}
    \item 早期错误假设可能在多轮会议中被反复引用，形成``共识幻觉''。
    \item 当所有角色共享同类训练偏差时，系统更易出现一致但错误的推理链。
    \item 如果没有显式反方角色，PI 的汇总常把``流畅度''误判为``正确性''。
\end{itemize}
\end{warningbox}

\subsection{实验结果与外部验证}

前一节说明 Critic 如何在候选进入实验前阻断过拟合与虚假高分；本节转向最关键的证据边界：计算候选是否真的在湿实验中显示预期结合信号。读曲线时应先看对照组、变体差异和重复趋势，而不是只看某一个最高点。

Virtual Lab 案例不是止于文字报告。课程展示了实验验证曲线和候选结构信息，说明至少在该任务上系统给出的候选在目标变体上表现出积极信号。图\ref{fig:vl-exp} 是相应关键帧。\footnote{来源：Berkeley RDI 课程视频，关键帧时间戳 00:18:54。}

\begin{figure}[H]
\centering
\includegraphics[width=0.92\textwidth]{frame-07-results-curve.jpg}
\caption{Virtual Lab 设计候选的实验验证示意（课程截图）}
\label{fig:vl-exp}
\end{figure}

\begin{knowledgebox}{读图：结构候选与实验曲线必须一起解释}
上半部分给出候选 nanobody 的结构与突变位置，下半部分比较不同抗原浓度下的结合信号。应关注新变体与原始株的曲线是否都高于未突变对照，以及趋势能否跨浓度保持。该图支持“候选值得进一步实验”，但不单独证明临床有效性、广谱性或长期安全性。
\end{knowledgebox}

要注意的是，Zou 的论证方式并非``AI 已经取代实验科学家''，而是``AI 团队可以更快生成值得实验资源投入的候选''。这点很关键，因为科研瓶颈常在实验预算和周期，任何能提高候选质量与优先级排序效率的系统，都能在真实研发中放大价值。

\begin{knowledgebox}{该案例可迁移的方法论}
\begin{itemize}
    \item 先建立可计算、可验证的候选生成循环，再追求更高复杂度目标。
    \item 在数据稀疏场景里引入保守角色（Critic）和外部验证关口。
    \item 把``模型输出''转成``实验可执行清单''作为交付物。
\end{itemize}
\end{knowledgebox}

\subsection{本章小结}

nanobody 案例给出的不是某个神奇 prompt，而是一套可以重复执行的科研闭环：多角色提出方案、Critic 抑制风险、实验侧验证候选、再反馈回下一轮迭代。

